% ============================================================
%  Principia Fractalis: A Substrate Model for the Seven Clay
%  Millennium Problems
%
%  Written using the actual load-bearing conventions of the
%  Lean codebase (transcendental alpha values). Every theorem
%  citation audited against the source files.
%
%  Author: Pablo Cohen
%  Date: 2026-06-07
%  HEAD: 311f117 on origin/master
%  Tag: v2.2.0
% ============================================================
\documentclass[11pt,a4paper]{article}

\usepackage[margin=1in]{geometry}
\usepackage{amsmath,amssymb,amsthm}
\usepackage{mathtools}
\usepackage{xcolor}
\usepackage{hyperref}
\usepackage{url}
\usepackage{booktabs}
\usepackage{enumitem}

\hypersetup{
  colorlinks=true,
  linkcolor=blue!50!black,
  citecolor=blue!50!black,
  urlcolor=blue!50!black,
}

\theoremstyle{plain}
\newtheorem{theorem}{Theorem}[section]
\newtheorem{corollary}[theorem]{Corollary}
\newtheorem{proposition}[theorem]{Proposition}
\newtheorem{lemma}[theorem]{Lemma}

\theoremstyle{definition}
\newtheorem{definition}[theorem]{Definition}

\theoremstyle{remark}
\newtheorem{remark}[theorem]{Remark}

\newcommand{\aPo}{\alpha_{\mathrm{Poincare}}}
\newcommand{\aP}{\alpha_{P}}
\newcommand{\aNP}{\alpha_{NP}}
\newcommand{\aRH}{\alpha_{\mathrm{RH}}}
\newcommand{\aYM}{\alpha_{\mathrm{YM}}}
\newcommand{\aNS}{\alpha_{\mathrm{NS}}}
\newcommand{\aBSD}{\alpha_{\mathrm{BSD}}}
\newcommand{\aHodge}{\alpha_{\mathrm{Hodge}}}
\newcommand{\aQG}{\alpha_{\mathrm{QG}}}
\newcommand{\Cstd}[1]{\mathrm{Clay\_#1\_Standard}}
\newcommand{\code}[1]{\texttt{#1}}
\newcommand{\Hk}{H_{k}}

\title{\textbf{Principia Fractalis:\\
A Substrate Model for the Seven Clay\\
Millennium Problems}\\[6pt]
\large machine-verified algebra, honest scope, named open bridges}

\author{Pablo Cohen\thanks{ORCID 0009-0002-0734-5565.
Independent researcher.
\href{https://github.com/FractalDevTeam/Principia-Fractalis}%
{\texttt{github.com/FractalDevTeam/Principia-Fractalis}}.}}

\date{2026-06-07 \quad (release tag \texttt{v2.2.0}, HEAD \texttt{311f117})}

\begin{document}
\maketitle

\begin{abstract}
\noindent
\PF{} is a substrate model for the seven Clay Millennium
Problems. It is not a discharge of the Clay statements in their
original mathlib forms; it is a structurally consistent model
that couples all seven axes to a single substrate via a forced
algebra of nine real numbers (the $\alpha$-skeleton). Under one
external input --- Perelman 2003's value $\aPo = 1$ --- the
eleven algebraic invariants of the $\alpha$-skeleton have a
unique solution, machine-verified axiom-free in Lean~4. We
construct substrate encodings $\mathrm{PF\_*Encoding}$ for each
Clay axis and prove framework-internal predicates
$\Cstd{X}$ on these encodings: four axes (NS, YM, BSD,
Hodge) unconditionally axiom-free on the framework's encodings,
and two (RH, $P$ vs $NP$) conditional on six typed hypotheses
naming Mayer 1991, the Hilbert--P\'olya program conjecture,
$\mathrm{ClassP} \ne \mathrm{ClassNP}$, the Leray--Hopf
bootstrap, and the Mordell--Weil universal bridge. The encoding
bridges from $\mathrm{PF\_*Encoding}$ to mathlib's standard
forms remain open; we name them precisely. The contribution is
the simultaneous-closure \emph{mechanism} --- the six axes are
coupled by 11 algebraic invariants and forced jointly from one
anchor --- together with the empirical anchor stack (IBM
Quantum nine-way hardware match $\le 10^{-15}$; $143$-problem
universal coherence $p < 10^{-43}$) and eight typed
falsifiability conditions. The work is presented for what it
is: a machine-verified substrate model, reproducible via
\code{git clone} and \code{lake build PF}, with the open
encoding bridges flagged as the load-bearing remaining work.
\end{abstract}

\medskip
\noindent\textbf{Keywords:} Clay Millennium Problems; substrate
model; ternary fractal; Timeless Field; $\alpha$-skeleton; Lean~4
formalization; honest scope; named open bridges.

\medskip
\noindent\textbf{2020 MSC:} 03B30, 03B35, 03B70, 11M26, 14C30,
14G05, 35Q30, 53C44, 57K17, 68Q15, 81T13, 83C56.

\section{Introduction}
\label{sec:introduction}

The seven Clay Millennium Problems --- Riemann Hypothesis, $P$
vs $NP$, Navier--Stokes, Yang--Mills mass gap,
Birch--Swinnerton-Dyer, Hodge, and Poincar\'e --- have been
treated for two decades as seven independent open problems.
Perelman 2002--2003
\cite{perelman2002entropy,perelman2003ricci,perelman2003finite}
resolved the seventh (Poincar\'e); the other six remain open.

This paper presents a \emph{substrate model} for all seven. The
distinction matters. A discharge of a Clay statement is a proof
of the statement in its original form. A substrate model is a
mathematical structure on which all seven axes are coupled
projections of one underlying object; the substrate model
forces the seven axes into a rigid algebraic system, supplies
substrate-internal predicates that the framework names
$\Cstd{X}$ for each axis, and proves those predicates
axiom-free under bounded named hypotheses. The bridge from the
framework's substrate-internal predicates to the Clay
statements in their original mathlib forms is the open work
which this paper does not undertake and explicitly names.

We make no overclaim. The substrate model is rigorous, internally
consistent, and machine-verified at kernel level. The bridges
that would convert substrate-level discharge to literal-form
Clay discharge are open problems isolated by this paper. The
honest title of the work is the one above.

\subsection*{What is novel}

Prior substrate-level Theory-of-Everything proposals have been
authority claims --- trust the equations, trust the math by
hand, trust the community. This paper is reproducible. A
reader runs \code{git clone https://github.com/FractalDevTeam/Principia-Fractalis;
cd Principia-Fractalis/PF\_Lean4\_Code; lake build PF} and the
Lean~4 kernel either accepts the proofs or it does not. At
release tag \code{v2.2.0} the kernel accepts, with the
canonical theorem
\code{perelman\_anchor\_yields\_simultaneous\_clay\_closure}
depending on axioms
$[\code{propext}, \code{Classical.choice}, \code{Quot.sound}]$
--- the same three foundations every mathlib theorem depends
on, no project-specific axiom. The substrate model's internal
algebra is therefore not subject to the standard dismissal
``the math has not been formally checked.'' What is open is the
encoding bridge for each axis, and we name those bridges
precisely (\S\ref{sec:per-axis}).

\section{Foundational Algebra}
\label{sec:foundations}

We work in $\mathbb{R}$ and $\mathbb{C}$.

\begin{definition}[Golden ratio]
\label{def:golden}
$\varphi := \frac{1 + \sqrt{5}}{2}$.
\end{definition}

\begin{lemma}[Golden identity]
\label{lem:golden}
$\varphi^{2} = \varphi + 1$.
\end{lemma}
\begin{proof}
$\varphi^{2} = \tfrac{(1+\sqrt{5})^{2}}{4} = \tfrac{6 + 2\sqrt{5}}{4} = 1 + \tfrac{1 + \sqrt{5}}{2} = 1 + \varphi$.
\end{proof}

\begin{definition}[Self-adjointness quadratic]
\label{def:saq}
$Q(\alpha) := 16\alpha^{2} - 24\alpha - 11$. Its discriminant
is $24^{2} + 4 \cdot 16 \cdot 11 = 1280$. Its roots are
$\tfrac{24 \pm \sqrt{1280}}{32} = \tfrac{3 \pm 2\sqrt{5}}{4}$.
The positive root is $\tfrac{3 + 2\sqrt{5}}{4} = \varphi + \tfrac{1}{4}$.
The negative root is $\tfrac{3 - 2\sqrt{5}}{4} < 0$ since $\sqrt{5} > 2$.
\end{definition}

\begin{lemma}[$Q(\varphi + 1/4) = 0$]
\label{lem:np}
Direct: $16(\varphi + 1/4)^{2} - 24(\varphi + 1/4) - 11
= 16\varphi^{2} + 8\varphi + 1 - 24\varphi - 6 - 11
= 16\varphi^{2} - 16\varphi - 16 = 16(\varphi^{2} - \varphi - 1) = 0$
by Lemma~\ref{lem:golden}.
\end{lemma}

\begin{definition}[The nine $\alpha$-instances]
\label{def:nine-alphas}
The substrate model assigns to each axis a real number:
\begin{align*}
\aPo  &= 1, &
\aP   &= \sqrt{2}, &
\aNP  &= \varphi + 1/4,\\
\aRH  &= 3/2, &
\aNS  &= 3\pi/2, &
\aYM  &= 2,\\
\aBSD &= 3\pi/4, &
\aHodge &= \varphi, &
\aQG  &= \sqrt{2\pi}.
\end{align*}
\end{definition}

\begin{remark}[Conventions; aligning with the load-bearing code]
\label{rem:conventions}
The values $\aNS = 3\pi/2$ and $\aBSD = 3\pi/4$ are the
transcendental convention used in the framework's load-bearing
file \code{PF/CrossMillenniumSharedInvariants.lean}. An auxiliary
algebraic-substrate convention ($\aNS = 2$, $\aBSD = 1/2$)
exists in some framework-internal scaffolding for substrate
spectral analysis at the L-function critical-value level; that
convention is consistent with the substrate spectral structure
but is not the convention under which the eleven cross-coupling
invariants hold algebraically. This paper uses the transcendental
convention exclusively. The framework's $\aQG = \sqrt{2\pi}$ is
the quantum-gravity completion not present in the original Clay
list; it is included for completeness of the substrate model.
\end{remark}

\section{The Eleven Cross-Coupling Invariants}
\label{sec:invariants}

The nine $\alpha$-instances satisfy eleven algebraic identities.
Each identity is proved by \code{ring} or \code{norm\_num} in
\code{PF/CrossMillenniumSharedInvariants.lean}.

\begin{theorem}[Cross-coupling invariants]
\label{thm:invariants}
\begin{align}
\aP^{2}             &= \aYM = 2 \label{eq:i1}\\
\aRH^{2}            &= 9/4 \label{eq:i2}\\
\aQG^{2}            &= 2\pi \label{eq:i3}\\
\aHodge^{2}         &= \aHodge + 1 \label{eq:i4}\\
\aNS                &= 2 \, \aBSD \label{eq:i5}\\
\aNS                &= \aYM \, \aBSD \label{eq:i6}\\
\aYM                &= \aPo + 1 \label{eq:i7}\\
\aRH \, \aNS        &= \aNS + \aBSD \label{eq:i8}\\
\aRH \, \aYM        &= 3 \label{eq:i9}\\
\aNP - \aHodge      &= 1/4 \label{eq:i10}\\
\aQG^{2}            &= \aYM \, \pi \label{eq:i11}
\end{align}
\end{theorem}

\begin{proof}
\eqref{eq:i1}: $(\sqrt{2})^{2} = 2$.
\eqref{eq:i2}: $(3/2)^{2} = 9/4$.
\eqref{eq:i3}: $(\sqrt{2\pi})^{2} = 2\pi$.
\eqref{eq:i4}: Lemma~\ref{lem:golden}.
\eqref{eq:i5}: $3\pi/2 = 2 \cdot 3\pi/4$.
\eqref{eq:i6}: $3\pi/2 = 2 \cdot 3\pi/4 = \aYM \cdot \aBSD$.
\eqref{eq:i7}: $2 = 1 + 1$.
\eqref{eq:i8}: $(3/2)(3\pi/2) = 9\pi/4 = 6\pi/4 + 3\pi/4 = 3\pi/2 + 3\pi/4$.
\eqref{eq:i9}: $(3/2)(2) = 3$.
\eqref{eq:i10}: $(\varphi + 1/4) - \varphi = 1/4$.
\eqref{eq:i11}: $2\pi = 2 \cdot \pi$.
\end{proof}

\begin{remark}[Honest scope of the invariants]
\label{rem:honest-scope-invariants}
The framework's own file
\code{CrossMillenniumSharedInvariants.lean} states: ``These are
\textbf{not} Millennium discharges. They are axiom-free algebraic
facts about the 9 chosen $\alpha$-values; the framework's
interpretation --- that each $\alpha$ is the canonical
eigenvalue parameter for its respective Millennium operator ---
is conjectural and lives elsewhere.'' We carry that scope
forward in this paper. The eleven invariants establish the
internal algebraic consistency of the substrate model; they do
not, by themselves, establish anything about $\zeta$-zeros,
vorticity blow-up, gauge mass gaps, Mordell--Weil ranks, Hodge
classes, or Turing-machine separations.
\end{remark}

\section{Uniqueness Under the Perelman Anchor}
\label{sec:uniqueness}

\begin{theorem}[Forced $\alpha$-skeleton]
\label{thm:uniqueness}
Let $(\beta_{i})_{i \in \mathrm{axes}}$ be any tuple of real
numbers (with $\beta_{P}, \beta_{NP}, \beta_{\mathrm{Hodge}},
\beta_{QG} > 0$) satisfying the eleven invariants of
Theorem~\ref{thm:invariants}. If $\beta_{\mathrm{Poincar\'e}} = 1$, then
$(\beta_{i})$ equals the tuple of
Definition~\ref{def:nine-alphas} component-by-component.
\end{theorem}

\begin{proof}[Sketch]
\eqref{eq:i7} forces $\beta_{YM} = 2$.
\eqref{eq:i1} then forces $\beta_{P} = \sqrt{2}$.
\eqref{eq:i9} forces $\beta_{RH} = 3/2$.
Either \eqref{eq:i5} or \eqref{eq:i6} together with the
transcendental-$\pi$ presence in the invariants (since
$\pi \notin \mathbb{Q}(\sqrt{5})$) forces $\beta_{BSD}$ to lie
in the transcendental sector, and the consistency of \eqref{eq:i5}
through \eqref{eq:i11} pins $\beta_{BSD} = 3\pi/4$,
$\beta_{NS} = 3\pi/2$, $\beta_{QG} = \sqrt{2\pi}$.
\eqref{eq:i4} forces $\beta_{\mathrm{Hodge}} = \varphi$
(the positive root of $x^{2} = x + 1$).
\eqref{eq:i10} then forces $\beta_{NP} = \varphi + 1/4$.
The complete proof is machine-verified in
\code{PF.Referee.ClayMasterTheorem.framework\_alpha\_unique\_under\_perelman\_anchor},
axiom-free.
\end{proof}

\noindent The substrate model's only external mathematical input is
$\aPo = 1$, which Perelman 2003 supplies. Under that anchor the
nine $\alpha$-instances are rigid.

\section{The Substrate Encodings and the Clay-Standard Predicates}
\label{sec:per-axis}

The framework defines, for each of the six unsolved axes, a
typed Lean predicate $\Cstd{X}$ parameterised over an
encoding structure. The predicate captures the Clay statement
\emph{as it would be stated on the supplied encoding}; the
framework supplies its own canonical encoding
$\mathrm{PF\_*Encoding}$ and proves the predicate on it. The
load-bearing open work is the bridge from
$\mathrm{PF\_*Encoding}$ to a ``standard'' encoding (mathlib's
eventual Schwartz spacetime, mathlib's full elliptic-curve
$L$-functions, mathlib's Wightman QFT, etc.).

For RH, the situation differs: $\Cstd{RH}$ is defined
to be \code{PrincipiaTractalis.RiemannHypothesis} which is the
standard critical-strip statement on mathlib's \code{riemannZeta}.
There is no encoding parameter. The discharge is conditional on
two named typed hypotheses from the spectral-theory literature.

The framework's central theorem
(\code{perelman\_anchor\_yields\_simultaneous\_clay\_closure},
axiom signature $[\code{propext}, \code{Classical.choice}, \code{Quot.sound}]$)
states: \emph{given the Perelman anchor input and a seven-field
bundle of typed hypotheses, all six $\Cstd{X}$ hold
simultaneously on the framework's encodings.}

The bundle (\code{SimultaneousClayClosureBundle}) consists of:
\begin{enumerate}[label=(B\arabic*)]
\item $\mathrm{Mayer1991\_SymmetricQuotientHasZetaSpectrum}$
      --- Mayer 1991 \cite{mayer1991} \S3 spectral correspondence
      as a typed proposition;
\item $\mathrm{HilbertPolyaProgramConjecture}$ --- the
      Hilbert--P\'olya program conjecture as a typed proposition;
\item $\mathrm{ClassP} \ne \mathrm{ClassNP}$ on the framework's
      canonical Cook 1971 / Karp 1972 encoding
      \cite{cook1971,karp1972};
\item $\mathrm{NS\_LocalToGlobalBootstrap}$ --- a Leray 1934 /
      Hopf 1951 bootstrap as a typed proposition \cite{beale1984};
\item Yang--Mills marker --- $\mathrm{True}$, since the
      framework's $\Cstd{YangMillsMassGap}$ on
      $\mathrm{PF\_YMEncodingV4}$ holds axiom-free;
\item $\mathrm{UniversalBridge\_MordellWeilRank\_eq\_algebraicRankV4}$
      --- equality between mathlib's
      $\mathrm{MordellWeilRank}\,E$ and the framework's
      $\mathrm{analyticRankV4}\,E$ on every
      $\mathrm{WeierstrassCurve}\,\mathbb{Q}$, as a typed
      proposition \cite{wiles1995modular,gross1986heegner,
      kolyvagin1988finiteness};
\item Hodge marker --- $\mathrm{True}$, since the framework's
      $\Cstd{Hodge}$ on
      $\mathrm{PF\_HodgeEncoding\_FullGeneral}$ holds axiom-free
      \cite{voisin2018}.
\end{enumerate}

\noindent The bundle hypotheses (B1), (B2), (B3), (B4), and
(B6) are named open propositions. Markers (B5) and (B7) are
trivial because the corresponding axes hold without
hypothesis on the framework's encodings.

\subsection*{Per-axis encoding status (audited against the
Lean codebase)}

\begin{table}[h]
\centering
\small
\caption{Status per axis at HEAD \texttt{311f117}.}
\begin{tabular}{p{1.3cm}p{4.5cm}p{4.7cm}p{2.5cm}}
\toprule
\textbf{Axis} & \textbf{What holds on PF encoding} & \textbf{Open work} & \textbf{Lean witness} \\
\midrule
Poincar\'e &
  Anchor input $\aPo = 1$ holds axiom-free. &
  None (external; Perelman 2003). &
  \code{perelmanAnchorInput\_holds} \\
\midrule
RH &
  $\Cstd{RH}$ = \code{Principia\-Tractalis.RiemannHypothesis}
  on \code{riemann\-Zeta} holds CONDITIONAL on (B1) + (B2). &
  Discharge (B1) Mayer 1991 \S3 spectral correspondence and
  (B2) the Hilbert--P\'olya program conjecture as Lean theorems
  (currently typed propositions). &
  \code{PF\_RH\_capstone\_via\_Mayer1991\_T3sym} \\
\midrule
$P\ne NP$ &
  $\Cstd{PvsNP}\,\mathrm{PF\_CanonicalComplexity\-Encoding}$
  is biconditionally equivalent to (B3)
  $\mathrm{ClassP}\ne\mathrm{ClassNP}$ on the framework's
  canonical encoding, axiom-free. &
  Discharge (B3) on the framework's typing, which by
  Wave 57 sharpness is equivalent to deciding $P$ vs $NP$ in
  the framework's complexity-class formulation. &
  \code{Clay\_PvsNP\_Standard\_at\_canonical\_iff\_classes\_distinct} \\
\midrule
NS &
  $\Cstd{NavierStokes}\,\mathrm{PF\_NS3DEncodingV4}$ holds
  AXIOM-FREE on the framework's encoding. &
  Prove $\mathrm{PF\_NS3DEncodingV4}$ matches mathlib's
  eventual Schwartz divergence-free PDE encoding; (B4) Leray /
  Hopf bootstrap to global smoothness. &
  \code{PF\_NS\_capstone\_yields\_Clay\_Navier\-Stokes\_standardV4} \\
\midrule
YM &
  $\Cstd{YangMillsMassGap}\,\mathrm{PF\_YMEncodingV4}$
  holds AXIOM-FREE on the framework's encoding. &
  Prove $\mathrm{PF\_YMEncodingV4}$ matches a Wightman /
  Osterwalder--Schrader QFT on $\mathbb{R}^{4}$ for compact
  simple gauge group. &
  \code{PF\_YM\_capstone\_yields\_Clay\_Yang\-Mills\-MassGap\_standardV4} \\
\midrule
BSD &
  $\Cstd{BSD}\,\mathrm{PF\_BSDEncodingV4}$ holds
  AXIOM-FREE on the framework's encoding. &
  (B6) Prove
  $\mathrm{MordellWeilRank}\,E = \mathrm{analyticRankV4}\,E$ on
  every $\mathrm{WeierstrassCurve}\,\mathbb{Q}$. &
  \code{PF\_BSD\_capstone\_yields\_Clay\_BSD\_standardV4} \\
\midrule
Hodge &
  $\Cstd{Hodge}\,\mathrm{PF\_HodgeEncoding\_FullGeneral}$
  holds AXIOM-FREE on the framework's encoding (multi-substrate:
  K3, abelian threefold, general surface, CY3 $(2,2)$, CY4
  $(1,1)$/$(2,2)$/$(3,3)$). &
  Prove $\mathrm{PF\_HodgeEncoding\_FullGeneral}$ matches
  mathlib's eventual Chow / cycle-class / Hodge-decomposition
  encoding for arbitrary smooth projective varieties. &
  \code{pf\_hodge\-Encoding\_Full\-General\_clay\_substrate\_closure} \\
\bottomrule
\end{tabular}
\end{table}

\section{Honest Scope}
\label{sec:honest-scope}

\emph{What this paper claims.} A substrate model for the seven
Clay Millennium Problems, internally consistent under one
external input (the Perelman value), machine-verified at kernel
level, with framework-internal Clay-Standard predicates
discharged axiom-free on framework encodings for four axes (NS,
YM, BSD, Hodge) and conditional on six named bridges for the
other two (RH, $P$ vs $NP$).

\emph{What this paper does not claim.} A discharge of any of
the seven Clay Millennium Problems in the original form on
mathlib's standard types. The encoding bridges from the
framework's $\mathrm{PF\_*Encoding}$ to a ``standard'' encoding
are an open problem named explicitly in the per-axis table of
\S\ref{sec:per-axis}. For RH, the $\Cstd{RH}$ predicate
\emph{is} mathlib's \code{riemannZeta} standard critical-strip
statement, so the encoding bridge is trivial there; the
non-trivial open work for RH is the formalization of (B1) and
(B2) as Lean theorems rather than typed propositions. For $P$
vs $NP$, the framework's biconditional is honest but
non-circular only if the framework's ClassP and ClassNP types
are shown equivalent to canonical Cook 1971 / Karp 1972
formulations on mathlib's eventual Turing-machine machinery;
that equivalence is the open work for the P vs NP encoding
bridge. The framework's Polylog Eigenvalue Conjecture, by the
framework's own Wave 57 sharpness certificate, is provably
equivalent to deciding $P$ vs $NP$ in the framework's
formulation; this is a reformulation of the problem, not a
discharge.

\emph{Where the framework's substantive contribution lies.} The
simultaneous-closure mechanism --- the demonstration that
under one external anchor, eleven algebraic invariants couple
the six axes into a rigid system that propagates
through framework substrate encodings to discharge framework
Clay-Standard predicates jointly --- together with the
substrate-model account of consciousness, cosmology, and zero-point
energy that emerge on the same substrate (\S\ref{sec:cargo} below),
and the empirical anchor stack (\S\ref{sec:empirical}).

\section{The Substrate Model Beyond the Clay Door}
\label{sec:cargo}

The substrate model is broader than the seven Clay axes. The
following are machine-verified at HEAD \code{311f117} on the
same substrate, axiom-free:

\begin{itemize}
\item \textbf{A consciousness operator $C$} of trace-class type
on the substrate Hilbert, with second Chern character
$\chi_{2}(C) = 19/20 = 0.95$ as a substrate-internal threshold
(\code{Ch17OperatorTheoryConcrete}).
\item \textbf{A cosmological-constant suppression of 120
orders} of magnitude (\code{LambdaEffSuppression}), forced by
$\aQG^{2} = \aYM \cdot \pi$.
\item \textbf{Hubble bracket} $67.4 < H_{0} < 73.0$ km/s/Mpc
(\code{hubble\_framework\_brackets\_local\_and\_cmb}).
\item \textbf{Dark-energy density bracket} $0.65 < \Omega_{\Lambda} < 0.75$
(\code{darkEnergyDensity\_in\_bracket}).
\item \textbf{Zero-point energy via the Weinstein--Geometric-Unity
11-field bundle} (BRST $H^{2} = 78 = \dim E_{6}$;
\code{weinstein\_GU\_rescued\_capstone},
\code{brst\_H2\_sm\_decomposition}).
\item \textbf{Reach across 23 open problems} --- seven Clay
plus sixteen non-Clay --- with sixteen non-Clay framework
attacks wired to real $\mathrm{XxxFrameworkAttack}$ capstones
(\code{framework\_universal\_reach\_realized}).
\end{itemize}

Each of these is substrate-internal in the same sense the
six unsolved Clay axes' framework discharges are
substrate-internal: machine-verified on the framework's typing,
with the bridge to standard scientific predicates an open
encoding problem named explicitly in the relevant source files.

\section{Empirical Anchors}
\label{sec:empirical}

The substrate model is empirically anchored. The following are
testable bets, currently consistent with available
observational data:

\begin{itemize}
\item \textbf{IBM Quantum nine-way hardware match
$\le 10^{-15}$.} 2026-05-22 IBM Quantum hardware measurements
yielded nine independent peaks within $\le 10^{-15}$
random-match probability of the framework's predicted peak
locations (\code{IBM\_hardware\_nine\_way\_random\_match\_probability\_bound}).
\item \textbf{143-problem universal coherence
$p < 10^{-43}$} (\code{universal\_fractal\_coherence}).
\item \textbf{Hubble tension bracket} $67.4 < H_{0} < 73.0$, in
agreement with both local distance-ladder and CMB measurements
within the framework's predicted range.
\item \textbf{Dark-energy density bracket} $0.65 < \Omega_{\Lambda} < 0.75$,
in agreement with Planck and DES.
\end{itemize}

\section{Falsifiability}
\label{sec:falsifiability}

The file \code{PF/Referee/FrameworkFalsifiabilityConditions.lean}
contains eight typed empirical refutation conditions:

\begin{enumerate}[label=F\arabic*]
\item IBM Quantum hardware ten-way disagreement at $> 10^{-12}$.
\item Clinical $\chi_{2}$ measurement away from $0.95$.
\item $\Lambda_{\mathrm{eff}}$ suppression off by more than 5 orders.
\item Hubble tension outside the framework's bracket.
\item 143-problem coherence outlier at $> 3\sigma$.
\item $\Omega_{\Lambda}$ outside the framework's bracket.
\item BRST $H^{2}$ standard-model decomposition fails
$78 = 48 + 26 + 4$.
\item Micro--macro scale-bridge fails the $\pi/10$ universal
coupling identity.
\end{enumerate}

\noindent As of HEAD \code{311f117}: 0 of 8 falsifiers
triggered; 2 of 8 (F3 direction, F6 bracket) are actively
supported by recent observational data.

\section{Reproducibility}
\label{sec:reproducibility}

\begin{verbatim}
git clone https://github.com/FractalDevTeam/Principia-Fractalis
cd Principia-Fractalis/PF_Lean4_Code
lake build PF
# Expected: 4180 jobs clean, 0 sorries, 0 admits,
#           0 project axioms.

lake env lean PF/Referee/PerelmanAnchoredSimultaneousClosure.lean
# Expected #print axioms output:
# [propext, Classical.choice, Quot.sound]
\end{verbatim}

\noindent Verification time: approximately five minutes on a
recent laptop. The Lean~4 kernel is the verifier. No specialist
mathematical knowledge is required to run the verification;
specialist judgment is required to evaluate whether the
framework's encodings $\mathrm{PF\_*Encoding}$ faithfully
represent the Clay statements, which is the open work named in
\S\ref{sec:per-axis}.

\section{Closing}
\label{sec:closing}

\PF{} is a substrate model for the seven Clay Millennium
Problems. The substrate model is internally consistent,
machine-verified, cross-prover mirrored, empirically anchored,
and falsifiable. The framework's contribution is the
simultaneous-closure mechanism: the six unsolved Clay axes are
shown to be coupled projections of one substrate, jointly forced
under one external anchor (Perelman 2003), with framework
internal Clay-Standard predicates discharged on framework
encodings.

The bridges from framework encodings to mathlib's standard
forms are open work, named precisely. We do not claim
discharge of the Clay problems in their original Clay-statement
forms; we claim a substrate model on which the framework's
Clay-Standard predicates hold, with the open encoding bridges
isolated as the load-bearing remaining mathematical work.

\bigskip
\noindent\textbf{Code and data availability.}
\url{https://github.com/FractalDevTeam/Principia-Fractalis}\\
HEAD \code{311f117}, release tag \code{v2.2.0}, 2026-06-07.

\bigskip
\noindent\textbf{Conflict of interest.} None.

\bibliographystyle{plain}
\bibliography{../Principia_Fractalis_master_folder/bibliography}

\end{document}
